%%%%%%%%%%  scan idx 113  |  folio 102  %%%%%%%%%%

\medskip
\noindent\textbf{Lema 4.7.}
La sucesión
\[
0\longrightarrow
H^0\!\left(
Y_0,
\underline{\operatorname{Hom}}_{\mathcal O_{Y_0}}
(\underline I/\underline I^2,\mathcal O_{Y_0})
\right)
\longrightarrow[\underline M_{X_0,Y_0}(k[\varepsilon])]
\longrightarrow[\underline M_{X_0}(k[\varepsilon])]
\]
es exacta, donde $\underline I$ es el haz de ideales que define
$Y_0$ en $X_0$; así,
\[
0\longrightarrow\underline I\longrightarrow\mathcal O_{X_0}
\longrightarrow\mathcal O_{Y_0}\longrightarrow0
\]
es exacta.

%%%%%%%%%%  scan idx 114  |  folio 103  %%%%%%%%%%

\medskip
\noindent\textit{Demostración.}
Sea $\underline B$ la subcategoría plena de
$\underline M_{X_0,Y_0}(k[\varepsilon])$ cuyos objetos tienen la
forma
\[
\bigl(k[\varepsilon];
X_0\times_{\operatorname{Spec}(k)}\operatorname{Spec}(k[\varepsilon]),
Y\bigr).
\]
Entonces
\[
[\underline B]
=\ker\!\left(
[\underline M_{X_0,Y_0}(k[\varepsilon])]
\longrightarrow[\underline M_{X_0}(k[\varepsilon])]
\right).
\]
La categoría $\underline B$ es el grupoide de las deformaciones de
$Y_0$ sobre $k[\varepsilon]$ dentro de
$X_0\times_{\operatorname{Spec}(k)}\operatorname{Spec}(k[\varepsilon])$.
Así, un objeto es un diagrama
conmutativo
\[
\begin{tikzcd}[column sep=huge,row sep=large]
Y \arrow[rr,"i_Y"] \arrow[dr,"\mathrm{flat}"']
&& X_0\times_{\operatorname{Spec}(k)}
   \operatorname{Spec}(k[\varepsilon]) \arrow[dl] \\
& \operatorname{Spec}(k[\varepsilon]) &
\end{tikzcd}
\]
cuya fibra especial es la inmersión cerrada fija $i_{Y_0}$. Una
flecha $(Y_1,i_{Y_1})\longrightarrow(Y_2,i_{Y_2})$ es un
$k[\varepsilon]$-morfismo $\Phi:Y_1\longrightarrow Y_2$
(necesariamente un isomorfismo) cuya fibra especial es
$\operatorname{id}_{Y_0}$ y para el cual
$i_{Y_2}\circ\Phi=i_{Y_1}$.

Equivalentemente, un objeto de $\underline B$ es un par $(Y,j_Y)$,
donde $Y$ es una deformación de $Y_0$ sobre $k[\varepsilon]$ y
$j_Y:Y\longrightarrow X_0$ es un $k$-morfismo cuyo compuesto
$Y_0\longrightarrow Y\xrightarrow{j_Y}X_0$ es $i_{Y_0}$. Una flecha
$(Y_1,j_{Y_1})\longrightarrow(Y_2,j_{Y_2})$ es un morfismo
$\Phi:Y_1\longrightarrow Y_2$ cuya fibra especial es
$\operatorname{id}_{Y_0}$ y para el cual
$j_{Y_2}\circ\Phi=j_{Y_1}$.

La asignación que envía $(Y,j_Y)$ al diagrama
\[
\begin{tikzcd}[column sep=large,row sep=large]
0 \arrow[r]
& \underline I \arrow[r] \arrow[d]
& \mathcal O_{X_0} \arrow[r] \arrow[d,"\sisi{\mathfrak{J}_Y}{j_Y}"]
& \mathcal O_{Y_0} \arrow[r] \arrow[d,equal]
& 0 \\
0 \arrow[r]
& \mathcal O_{Y_0} \arrow[r,"\varepsilon"]
& \mathcal O_Y \arrow[r]
& \mathcal O_{Y_0} \arrow[r]
& 0
\end{tikzcd}
\]
da un isomorfismo
\[
\begin{aligned}
[\underline B]
&\xrightarrow{\ j_Y|_{\underline I}\ }
\operatorname{Hom}_{\mathcal O_{X_0}}
 (\underline I,\mathcal O_{Y_0}) \\
&=
\operatorname{Hom}_{\mathcal O_{Y_0}}
 (\underline I/\underline I^2,\mathcal O_{Y_0}) \\
&=
H^0\!\left(
Y_0,
\underline{\operatorname{Hom}}_{\mathcal O_{Y_0}}
(\underline I/\underline I^2,\mathcal O_{Y_0})
\right).
\end{aligned}
\]
Esto demuestra la exactitud. \hfill$\square$

%%%%%%%%%%  scan idx 115  |  folio 104  %%%%%%%%%%

\medskip
\noindent\textbf{Observación 4.8.}
Para una fórmula más general que implica el Lema 4.7, véase SGA 3,
III, 4.4. También puede observarse lo siguiente. Para todo
$S$-esquema $X$ y todo $\mathcal O_X$-módulo cuasicoherente $E$, se
pueden definir $\mathcal O_X$-módulos cuasicoherentes
$\underline D^i_{X|S}(E)$ para todo $i\geq0$, generalizando el caso
en que $S$ es afín. Como de costumbre, escribimos
\[
\underline D^i_{X|S}(\mathcal O_X)=\underline D^i_{X|S}.
\]
Para la inmersión cerrada $Y_0\hookrightarrow X_0$, tenemos entonces
\[
H^0\!\left(
Y_0,
\underline{\operatorname{Hom}}_{\mathcal O_{Y_0}}
(\underline I/\underline I^2,\mathcal O_{Y_0})
\right)
=H^0(Y_0,\underline D^1_{Y_0|X_0}),
\]
donde
$0\longrightarrow\underline I\longrightarrow\mathcal O_{X_0}
\longrightarrow\mathcal O_{Y_0}\longrightarrow0$ es exacta.

\medskip
\noindent\textbf{Teorema 4.9.}
Sea $X_0$ un $k$-esquema de tipo finito y $Y_0$ un subesquema
cerrado de $X_0$. Supongamos que se cumple una de las condiciones
siguientes:

\smallskip
\noindent(i) $X_0$ es propio sobre $k$, o bien

\noindent(ii) $X_0$ es afín y liso sobre $k$ fuera de un número
finito de puntos, y $Y_0$ es propio sobre $k$.

\smallskip
\noindent\textbf{(a)}
La categoría $\underline M_{X_0,Y_0}$ admite una deformación formal
versal sobre $\underline C_\Lambda$. Prorepresenta el funtor
$[\underline M_{X_0,Y_0}]$ si y solo si, para toda aplicación
sobreyectiva $R'\longrightarrow R$ en $\underline C_\Lambda$ y toda
deformación $(R';X',Y')$, la aplicación
\[
\operatorname{Aut}_{R'}((X',Y')\mid X_0)
\longrightarrow
\operatorname{Aut}_{R}\!\left(
(X'\otimes_{R'}R,Y'\otimes_{R'}R)\mid X_0
\right)
\]
es sobreyectiva.

\smallskip
\noindent\textbf{(b)}
Si, además,
\[
H^0\!\left(
X_0,
\underline D^0_{Y_0|k}
\times_{\underline D^0_{X_0|k}(\mathcal O_{Y_0})}
\underline D^0_{X_0|k}
\right)
\]
tiene dimensión finita, entonces $\underline M_{X_0,Y_0}$ es
pseudorrepresentable de manera lisa.

\medskip
\noindent\textit{Demostración.}
Esto se deduce inmediatamente de los Lemas 4.4 y 4.7, junto con 1.11,
2.7 y 2.13. \hfill$\square$
